% META: {"id": "1SPE-TRIGO-EX-023", "chapitre": "1SPE-TRIGONOMETRIE", "type_objet": "exercice", "capacites": ["1SPE-TRIGONOMETRIE-C3"], "capacites_codes": ["C3"], "methodes": ["M3"], "parcours": 1, "competences": ["calculer"], "duree_min": 8, "mode_creation": "ex_nihilo", "sources_inspiration": [], "parametres_sympy": {}, "fichier_tex": "chapitres/1SPE-TRIGONOMETRIE/exercices/1SPE-TRIGO-EX-023.tex", "corrige_tex": "chapitres/1SPE-TRIGONOMETRIE/corriges/1SPE-TRIGO-CO-023.tex", "status": "generated"}
% BEGIN-VERIFY
% from sympy import *
% # cos(2*pi/3) using duplication: cos(2*pi/3) = 2*cos^2(pi/3) - 1 = 2*(1/4) - 1 = -1/2
% assert 2*Rational(1,4) - 1 == Rational(-1,2)
% assert cos(2*pi/3) == Rational(-1,2)
% # sin(2*pi/3) = 2*sin(pi/3)*cos(pi/3) = 2*(sqrt(3)/2)*(1/2) = sqrt(3)/2
% assert sin(2*pi/3) == sqrt(3)/2
% # sin(pi/3) = sin(2*pi/6) = 2*sin(pi/6)*cos(pi/6) = 2*(1/2)*(sqrt(3)/2) = sqrt(3)/2
% assert sin(pi/3) == sqrt(3)/2
% END-VERIFY
\begin{exercice}{1SPE-TRIGO-EX-023}{1}{8}
A l'aide des formules de duplication, calculer :
\begin{enumerate}
  \item $\cos\!\left(\dfrac{2\pi}{3}\right)$ en utilisant $\dfrac{2\pi}{3} = 2 \times \dfrac{\pi}{3}$.
  \item $\sin\!\left(\dfrac{2\pi}{3}\right)$ en utilisant $\sin(2a) = 2\sin a \cos a$.
  \item Retrouver $\sin\!\left(\dfrac{\pi}{3}\right)$ en utilisant $\dfrac{\pi}{3} = 2 \times \dfrac{\pi}{6}$.
\end{enumerate}
\end{exercice}
